\documentclass[10pt,a4paper]{article}
\usepackage[a4paper,margin=14mm,top=12mm,bottom=14mm]{geometry}
\usepackage{fontspec}
\setmainfont{Liberation Sans}
\usepackage{polyglossia}
\setdefaultlanguage{french}
\usepackage{amsmath}
\usepackage{xcolor}
\usepackage{graphicx}
\usepackage[most]{tcolorbox}
\usepackage{tabularx,booktabs,array,multirow}
\usepackage{enumitem}
\usepackage{fancyhdr}
\usepackage{lastpage}
\usepackage[hidelinks]{hyperref}
\definecolor{uimm}{HTML}{005677}
\definecolor{cyan}{HTML}{13A2DD}
\definecolor{fond}{HTML}{E6F2FA}
\definecolor{ok}{HTML}{2E7D32}
\definecolor{ko}{HTML}{B3261E}
\definecolor{violet}{HTML}{7B1FA2}
\setlength{\parindent}{0pt}
\setlength{\parskip}{3pt}
\setlist{nosep}
\setlist[itemize]{leftmargin=13pt,label=\textcolor{cyan}{\textbullet}}
\newcommand{\titre}[1]{\vspace{5pt}{\color{uimm}\large\bfseries #1}\par\vspace{-3pt}{\color{cyan}\rule{\linewidth}{1.2pt}}\par\vspace{0pt}}
\tcbset{boite/.style={colback=fond,colframe=uimm,boxrule=0.6pt,arc=2pt,left=6pt,right=6pt,top=4pt,bottom=4pt},
        defi/.style={colback=violet!6,colframe=violet,boxrule=0.8pt,arc=2pt,left=6pt,right=6pt,top=4pt,bottom=4pt,
                     title={\bfseries Défi},coltitle=white,colbacktitle=violet,fonttitle=\small}}
\newcommand{\bandeau}[3]{%
\begin{tcolorbox}[colback=uimm,colframe=uimm,arc=2pt,left=4pt,right=8pt,top=4pt,bottom=4pt]
\begin{minipage}[c]{18mm}\setlength{\fboxsep}{1.2pt}\fcolorbox{white}{white}{\includegraphics[height=16.5mm]{logo_uimm.jpg}}\end{minipage}\hspace{4mm}%
\begin{minipage}[c]{\dimexpr\linewidth-22mm\relax}\color{white}
{\small #1}\\[1pt]{\LARGE\bfseries #2}\\[2pt]{\large #3}\hfill{\small CQPM Technicien qualité}
\end{minipage}
\end{tcolorbox}}
\newcommand{\reponse}[1][1]{\par\vspace{2pt}{\color{gray!60}\foreach \i in {1,...,#1}{\rule{\linewidth}{0.3pt}\par\vspace{9pt}}}}
\usepackage{pgffor}
\pagestyle{fancy}\fancyhf{}
\renewcommand{\headrulewidth}{0pt}
\newcommand{\pied}[1]{\fancyfoot[L]{\footnotesize\color{uimm}#1}\fancyfoot[R]{\footnotesize\color{uimm}Page \thepage{} sur \pageref{LastPage}}}

\usepackage{tikz}
\usepackage{longtable}
\renewcommand{\bandeau}[3]{%
\begin{tcolorbox}[colback=uimm,colframe=uimm,arc=2pt,left=4pt,right=8pt,top=4pt,bottom=4pt]
\begin{minipage}[c]{18mm}\setlength{\fboxsep}{1.2pt}\fcolorbox{white}{white}{\includegraphics[height=16.5mm]{logo_uimm.jpg}}\end{minipage}\hspace{4mm}%
\begin{minipage}[c]{\dimexpr\linewidth-22mm\relax}\color{white}
{\small #1}\\[1pt]{\LARGE\bfseries #2}\\[2pt]{\large #3}
\end{minipage}
\end{tcolorbox}}
% en-tête d'une fiche : code, titre, question, bloc, durée, organisation, niveau
\newcommand{\entete}[7]{%
\begin{tcolorbox}[colback=uimm,colframe=uimm,arc=2pt,left=4pt,right=6pt,top=4pt,bottom=4pt]
\begin{minipage}[c]{18mm}\setlength{\fboxsep}{1.2pt}\fcolorbox{white}{white}{\includegraphics[height=16.5mm]{logo_uimm.jpg}}\end{minipage}\hspace{3mm}%
\begin{minipage}[c]{\dimexpr\linewidth-58mm\relax}\color{white}\raggedright
{\small Base SPC · #4}\\[1pt]{\Large\bfseries #1 · #2}\\[2pt]{\itshape #3}
\end{minipage}\hfill
\begin{minipage}[c]{33mm}\color{white}\small\raggedleft
Durée : \textbf{#5}\\ \textbf{#6}\\ Niveau #7 sur 3
\end{minipage}
\end{tcolorbox}}
\newcommand{\tagE}{{\setlength{\fboxsep}{1.5pt}\colorbox{cyan}{\color{white}\scriptsize\bfseries Ellistat}}}
\newcommand{\tagW}{{\setlength{\fboxsep}{1.5pt}\colorbox{uimm}{\color{white}\scriptsize\bfseries Page web}}}
\tcbset{pred/.style={colback=orange!6,colframe=orange!80!black,boxrule=0.6pt,arc=2pt,left=6pt,right=6pt,top=3pt,bottom=1pt,
                    title={\bfseries Je prédis, avant de commencer},coltitle=white,colbacktitle=orange!80!black,fonttitle=\small},
        retenir/.style={colback=fond,colframe=uimm,boxrule=0.6pt,arc=2pt,left=6pt,right=6pt,top=3pt,bottom=3pt,
                    title={\bfseries À retenir},coltitle=white,colbacktitle=uimm,fonttitle=\small},
        plus/.style={colback=violet!6,colframe=violet,boxrule=0.6pt,arc=2pt,left=6pt,right=6pt,top=3pt,bottom=3pt,
                    title={\bfseries Pour aller plus loin},coltitle=white,colbacktitle=violet,fonttitle=\small},
        act/.style={colback=uimm!10,colframe=uimm,boxrule=0.5pt,arc=2pt,left=5pt,right=5pt,top=2pt,bottom=2pt,before skip=7pt,after skip=3pt},
        verif/.style={colback=orange!6,colframe=orange!80!black,boxrule=0.6pt,arc=2pt,left=6pt,right=6pt,top=3pt,bottom=3pt,
                    title={\bfseries À vérifier sur poste avant la première séance},coltitle=white,colbacktitle=orange!80!black,fonttitle=\small}}
\newcommand{\fichier}[1]{\texttt{\small #1}}
\newcommand{\menu}[1]{\textsf{\textbf{#1}}}
\newcommand{\champ}[1]{\rule{#1}{0.3pt}}

\usetikzlibrary{decorations.pathreplacing}
\usepackage{colortbl}
\pied{Base SPC · Mémo · Le Gage R\&R}
\newcommand{\stitre}[1]{\par\vspace{1pt}{\color{uimm}\bfseries\large #1}\par\vspace{-7pt}{\color{cyan}\rule{\linewidth}{0.9pt}}\par\vspace{-4pt}}
\begin{document}
\thispagestyle{fancy}
\fontsize{9.3}{11}\selectfont
\bandeau{Base SPC · Mémo}{Le Gage R\&R, c'est quoi ?}{Peut-on se fier à la mesure avant de juger la pièce ?}
\vspace{2pt}
\begin{tcolorbox}[boite,top=2pt,bottom=2pt]
\textbf{L'idée.} Avant de dire qu'une pièce est bonne ou mauvaise, on vérifie que la \textbf{mesure} est fiable. Un \textbf{Gage R\&R} fait cette vérification. \emph{Gage} veut dire instrument de mesure. \emph{R\&R} veut dire \textbf{R}épétabilité et \textbf{R}eproductibilité. Le résultat se donne en pourcentage de la tolérance : le \textbf{\%GRR}. Plus il est petit, plus la mesure est fiable.
\end{tcolorbox}

\stitre{1. D'où vient la variation que l'on observe ?}
Dix pièces donnent dix valeurs différentes : les \textbf{pièces sont vraiment différentes} (on veut le voir) et la \textbf{mesure n'est pas parfaite} (le parasite). Le Gage R\&R chiffre le parasite, qui a deux composantes.\par\vspace{2pt}
\renewcommand{\arraystretch}{1.06}
\begin{tabularx}{\linewidth}{|>{\bfseries\raggedright\arraybackslash}p{24mm}|>{\raggedright\arraybackslash}X|>{\raggedright\arraybackslash}X|}\hline
\rowcolor{uimm}& \color{white}\textbf{Répétabilité (EV) : moi, je répète} & \color{white}\textbf{Reproductibilité (AV) : un autre reproduit}\\\hline
La question & Une même personne, avec le même instrument, remesure la même pièce. Trouve-t-elle les mêmes valeurs ? & Des personnes différentes, avec le même instrument, mesurent la même pièce. Trouvent-elles la même moyenne ?\\\hline
Exemple imaginé & Jean trouve 11,92 ; 11,91 ; 11,92. Ses valeurs sont groupées. & Jean trouve environ 11,92 et Marie environ 11,95. Chacun est régulier, mais ils ne sont pas d'accord.\\\hline
Comment on la chiffre & L'écart-type des mesures répétées (même personne, même pièce). & L'écart-type entre les moyennes des personnes.\\\hline
\end{tabularx}

\stitre{2. Comment se calcule le GRR ?}
\begin{minipage}[c]{0.30\linewidth}
\centering
\begin{tikzpicture}[x=0.62cm,y=0.62cm,font=\small]
\fill[fond] (0,0) -- (3,0) -- (0,4) -- cycle;
\draw[orange!85!black,line width=1.6pt] (0,0) -- (0,4) node[midway,left=1pt,black]{EV\,=\,0,004};
\draw[red!70!black,line width=1.6pt] (0,0) -- (3,0) node[midway,below=1pt,black]{AV\,=\,0,003};
\draw[uimm,line width=1.6pt] (3,0) -- (0,4) node[midway,above right=0pt,black]{\bfseries GRR};
\draw (0,0.35) -- (0.35,0.35) -- (0.35,0);
\end{tikzpicture}
\end{minipage}\hfill
\begin{minipage}[c]{0.67\linewidth}
Les deux erreurs ne s'additionnent pas : elles se composent comme les côtés d'un \textbf{triangle rectangle}, car ce sont leurs carrés (les variances) qui s'additionnent.
\begin{center}$\mathrm{GRR}=\sqrt{\mathrm{EV}^2+\mathrm{AV}^2}$\end{center}\vspace{-4pt}
Exemple : $\sqrt{0{,}004^2+0{,}003^2}=\sqrt{0{,}000025}=0{,}005$ mm. Le GRR est plus petit que $\mathrm{EV}+\mathrm{AV}=0{,}007$. Ellistat calcule EV, AV et GRR (par analyse de variance) : vous les lisez à l'écran en G2, sans calcul à la main.
\end{minipage}

\stitre{3. Dans Ellistat, « Gage R\&R » est d'abord le nom d'un menu}
\renewcommand{\arraystretch}{1.06}
\begin{tabularx}{\linewidth}{|>{\bfseries}p{20mm}|>{\raggedright\arraybackslash}X|>{\raggedright\arraybackslash}p{28mm}|}\hline
\rowcolor{uimm}\color{white}\textbf{Onglet} & \color{white}\textbf{Question posée} & \color{white}\textbf{Fiches}\\\hline
Justesse & Le micromètre est-il juste sur une cale connue ? & M3, M7, M16\\\hline
Linéarité & Est-il juste sur toute la longueur mesurée ? & M8, G1\\\hline
\rowcolor{orange!12}Mesures & Les opérateurs et l'instrument, ensemble, sont-ils fiables ? \textbf{C'est le R\&R.} & G2\\\hline
Attributs & Trois contrôleurs qui trient « bon » ou « rejet » décident-ils pareil ? & G4\\\hline
\end{tabularx}
\par\vspace{2pt}
M3 est donc une étude de \textbf{justesse}, pas un R\&R. Le vrai R\&R se fait en G2 : chaque stagiaire est opérateur.

\stitre{4. Lire le résultat}
\begin{minipage}[t]{0.47\linewidth}\vspace{0pt}
\begin{tcolorbox}[boite,top=1pt,bottom=2pt]
\textbf{Une seule formule}\begin{center}\large $\%\mathrm{GRR}=\dfrac{6\times \mathrm{GRR}}{T}$\end{center}\vspace{-6pt}
GRR et $T$ (la tolérance) : même unité, ici le mm. Le 6 : la mesure s'étale sur environ 6 écarts types ($\pm 3$). On compare cette largeur à la tolérance.
\end{tcolorbox}
\renewcommand{\arraystretch}{1.1}
\begin{tabularx}{\linewidth}{|>{\raggedright\arraybackslash}p{26mm}|>{\raggedright\arraybackslash}X|}\hline
\cellcolor{ok!25}moins de 10 \% & capable\\\hline
\cellcolor{orange!25}10 à 30 \% & capable sous conditions\\\hline
\cellcolor{ko!25}plus de 30 \% & non capable\\\hline
\end{tabularx}
\par\smallskip
Il faut aussi au moins \textbf{5 catégories distinctes} (groupes de pièces que la mesure sait séparer).
\end{minipage}\hfill
\begin{minipage}[t]{0.50\linewidth}\vspace{0pt}
\textbf{Exemple.} Suite du calcul : GRR = 0,005 mm.\par\smallskip
\renewcommand{\arraystretch}{1.15}
\begin{tabularx}{\linewidth}{|c|>{\raggedright\arraybackslash}X|c|}\hline
\rowcolor{uimm}\color{white}\textbf{Tolérance} & \color{white}\textbf{Calcul} & \color{white}\textbf{\%GRR}\\\hline
0,4 mm & $6\times 0{,}005 / 0{,}4$ & \cellcolor{ok!25}7,5 \%\\\hline
0,08 mm & $6\times 0{,}005 / 0{,}08$ & \cellcolor{ko!25}37,5 \%\\\hline
\end{tabularx}
\par\smallskip
Mêmes mesures : seule la tolérance change, et le verdict passe de capable à non capable. Une mesure est bonne ou mauvaise \textbf{pour une tolérance donnée}.
\end{minipage}

\stitre{5. Pour vérifier que vous avez compris}
\begin{enumerate}[leftmargin=16pt,itemsep=1pt]
\item Même butée, même micromètre. Paul trouve 11,90 ; 11,91 ; 11,90. Léa trouve 11,95 ; 11,96 ; 11,95. Quel est le problème : la répétabilité ou la reproductibilité ?
\item EV $= 0{,}006$ mm et AV $= 0{,}008$ mm. Quel est le GRR ? Quel est le \%GRR pour $T = 0{,}4$ mm ? Quel verdict ?
\item Vous êtes dans Ellistat, menu « Gage R\&R », onglet Justesse. Faites-vous un R\&R ?
\end{enumerate}
{\footnotesize\color{gray}\textbf{Réponses.} 1. La reproductibilité : chacun est régulier, mais ils ne trouvent pas la même moyenne. 2. GRR $=\sqrt{0{,}006^2+0{,}008^2}=0{,}010$ mm (et non 0,014). \%GRR $=6\times 0{,}010/0{,}4=15$ \% : capable sous conditions. 3. Non : c'est une étude de justesse. Le R\&R se fait dans l'onglet Mesures.}

{\scriptsize\color{gray}\textbf{Sources.} AIAG, \emph{Measurement Systems Analysis}, 4\textsuperscript{e} éd. ; Bosch, \emph{Technical Statistics}, Booklet 10. Valeurs des exemples imaginées.}
\end{document}
